\section{Discussion and extensions}
\label{s:extensions}

This analysis achieves 3 of the 4 purposes it set out in section \ref{s:intro} well.
I develop robust data science techniques appropriate for the task
that allow efficient handling of mass spatial data of any size required
for a credible spatial analysis, for which Fame likely represents the upper bound of observations.
I apply the econometric theory of deploying higher-order spatial lags
in a contemporaneous peer effect model,
and demonstrates that this tool is useful in the critical
setting of analysing spatial productivity.
In this, I offer key insights on the contribution and importance
of different specifications varying weights of spatial peers,
and on the employment of different weighting structures, such as multiple non-overlapping groups,
to improve identification.
\par
I extend this to the dynamic setting: I show that alternative deeper lags,
this time on the temporal dimension, are permissible, aligning
theory presented on dynamic panel structures from \cite{Arellano1991}
and estimates proposed by \cite{Blundell1998} to a peer effects application,
again, particularly that the policy-crucial question of spatial productivity
can be investigated through these tools.
I offer insights on the complimentary and overlapping nature
of spatial and temporal lags, establishing a framework of investigation
to evaluate the performance of selection of instruments.
Additionally, I find some early evidence supporting the use of a static model
in this question, with minimal significant results for dynamic effects
for the spatial productivity spillover.
Taken seriously, this provokes re-examination of the centrifugal forces
occurring over time, inspiring colocation and urban development,
if it implies that firms believe the benefits of positive productivity shocks
to be durable.
\par
Nevertheless, the last purpose listed, aiming to achieve robust empirical
estimates of the spillover are inconsistent.
\itx{A priori} expectations of the key parameters estimated, $\delta_0,
\xi_0, \zeta_0$ would expect them to be positively signed 
A negative sign on some taken as truth, would imply competing,
centrifugal effects such as congestion or constrained input resources.
Even then, results presented are unusual for finding negative signs
on some coefficients and not others,
indeed specifications seem to consistently imply that
$\delta_0 > 0 \implies \xi_0<0, \zeta_0<0$, and vice versa.
This is an unusual pattern that most likely suggests misspecification:
models are most likely trying to fit around some combination
of the parameters of interest where the factors compete for contribution,
and a sudden attribution to one via random noise signifies a decrease to the other.
Alternatively, the unusual pattern of parameter movement
could be some accounting 
\par
This paper handles several complicated econometric issues that require
significant further scrutiny to yield a definitive answer to the nature
of these results. Any one of the uncertain model choices discussed in this paper,
likely in order of significance, the arbitrary deselection of peers
by a 5km distance boundary, the choice to only construct distance
between peers by spatial distance, could have been the key to invalid peer identity.
Further econometric investigation is required those which are most key.
From this, robust, policy-relevant econometric estimates can be derived.